\section*{Chapter \ref{ch:g9:randomness-and-simulation} --- Randomness and Simulation}

\begin{solution}{exo:g9:randomness-and-simulation:1}
Yes: the same \omterm{def:g9:randomness-and-simulation:seed}{seed} gives the same sequence, so both print the ten numbers of
the first line of \cref{out:g9:randomness-and-simulation:seed}. With \omterm{def:g9:randomness-and-simulation:seed}{seed} 2027,
the sequence is a different one.
\end{solution}

\begin{solution}{exo:g9:randomness-and-simulation:2}
$0 \to 1 \to 4 \to 5 \to 0$: $3 \times 5 + 1 = 16$, whose remainder by 8 is 0.
The cycle has 4 numbers only; 2, 3, 6 and 7 never appear from this \omterm{def:g9:randomness-and-simulation:seed}{seed}. A
generator with a short cycle is a poor one.
\end{solution}

\begin{solution}{exo:g9:randomness-and-simulation:3}
\code{"red"}, \code{"blue"} or \code{"green"}: one \omterm{def:g7:lists:list}{item} of the \omterm{def:g7:lists:list}{list}, each \omterm{def:g7:lists:list}{item}
with the same chance. Red fills 3 of the 5 places, so it comes out with
probability $3/5$.
\end{solution}

\begin{solution}{exo:g9:randomness-and-simulation:4}
\code{run} counts the heads in the current run and goes back to 0 at each tail;
\code{longest} keeps the best run so far:

\omcode{code/grade-9/05-randomness-and-simulation/exercises/longest_heads.py}{1}{13}{The
longest run of heads.}\label{lst:g9:randomness-and-simulation:sol-heads}

With \omterm{def:g9:randomness-and-simulation:seed}{seed} 9, it prints 7. Runs of 6 or 7 heads in 100 tosses are common, which
surprises most people.
\end{solution}

\begin{solution}{exo:g9:randomness-and-simulation:5}
Seven counts, each close to 10,000: the run with \omterm{def:g9:randomness-and-simulation:seed}{seed} 3 prints
\code{[10017, 9979, 10175, 9874, 9984, 9905, 10066]}. \code{randint(0, 6)}
includes both ends, so this ``die'' has seven faces, 0 to 6, each with
probability $1/7$ instead of $1/6$; a six comes up less often than it should.
\end{solution}

\begin{solution}{exo:g9:randomness-and-simulation:6}
\code{random.seed(2026)} runs at every turn of the \omterm{def:g4:algorithms-repeat-and-decide:loop}{loop}, so the generator
starts again from the same \omterm{def:g9:randomness-and-simulation:seed}{seed} each time and gives its first number, 1, every
time. Set the \omterm{def:g9:randomness-and-simulation:seed}{seed} once, before the \omterm{def:g4:algorithms-repeat-and-decide:loop}{loop}:

\omcode{code/grade-9/05-randomness-and-simulation/exercises/seed_loop_fixed.py}{1}{7}{The
\omterm{def:g9:randomness-and-simulation:seed}{seed} set once.}\label{lst:g9:randomness-and-simulation:sol-seed}

It prints \code{[1, 3, 5, 5, 6, 1, 2, 5]}, the start of the first line of
\cref{out:g9:randomness-and-simulation:seed}.
\end{solution}

\begin{solution}{exo:g9:randomness-and-simulation:7}
Only the ranks matter, so the deck holds each rank four times:

\omcode{code/grade-9/05-randomness-and-simulation/exercises/pair_chance.py}{1}{15}{Two
cards of the same rank.}\label{lst:g9:randomness-and-simulation:sol-pair}

It prints \code{simulated: 0.0587   exact: 3/51 = 0.0588}. Exactly: whatever the
first card, 3 of the 51 cards left have its rank.
\end{solution}

\begin{solution}{exo:g9:randomness-and-simulation:8}
\omcode{code/grade-9/05-randomness-and-simulation/exercises/grid_walk.py}{1}{16}{A
walk on a \omterm{def:g9:grids-and-images:grid}{grid}.}\label{lst:g9:randomness-and-simulation:sol-grid}

With \omterm{def:g9:randomness-and-simulation:seed}{seed} 6, the token ends at \code{2 4}, 6 steps from home along the \omterm{def:g9:grids-and-images:grid}{grid}:
2 across and 4 up.
\end{solution}

\begin{solution}{exo:g9:randomness-and-simulation:9}
Of the 36 pairs of faces, 6 are doubles, $(1, 1)$ to $(6, 6)$: $6/36 = 1/6$.
\end{solution}

\begin{solution}{exo:g9:randomness-and-simulation:10}
Every number it gives is a remainder by 16, so one of the 16 \omterm{def:g6:variables-and-input:value}{values} 0 to 15.
Among any 17 numbers in a row, two must be equal. And once a \omterm{def:g6:variables-and-input:value}{value} comes back,
the rule computes the same next \omterm{def:g6:variables-and-input:value}{value} as the first time, and the next, and so
on: the sequence goes round the same cycle forever. A cycle of 16 is the most
it can do; the generator of the chapter reaches it.
\end{solution}

\begin{solution}{exo:g9:randomness-and-simulation:11}
In \cref{out:g9:randomness-and-simulation:frequencies}, the frequencies for 10
to 500 throws range from 0.10 to 0.20; from 10,000 throws on, they stay between
0.1657 and 0.1683. With 100 throws, one six more or less moves the frequency by
0.01; with 100,000 throws, by 0.00001: a few lucky or unlucky throws cannot
move it far. With 60 throws, each six moves the frequency by about 0.017, so 4
sixes instead of 10 is a big gap, but it is made of only 6 throws: Tom should
throw again, many more times, before blaming the die.
\end{solution}

\begin{solution}{exo:g9:randomness-and-simulation:12}
The \omterm{def:g8:dictionaries:dictionary}{dictionary} \code{ways} holds, for each position, the number of paths of the
steps made so far that reach it; each step sends every path one place left and
one place right:

\omcode{code/grade-9/05-randomness-and-simulation/exercises/walk_return.py}{1}{21}{Back
home after 20 steps.}\label{lst:g9:randomness-and-simulation:sol-return}

It prints \code{simulated: 0.1768}, then \code{exact: 184756 paths of 1048576,
0.1762}: of the $2^{20}$ equally likely walks, 184,756 end at 0. The \omterm{def:g9:randomness-and-simulation:simulation}{simulation}
is right to about a thousandth.
\end{solution}

\begin{solution}{pb:g9:randomness-and-simulation:1}
\textbf{1.} There are only 365 possible birthdays: once 365 pupils all have
different ones, the 366th must share one of them.

\textbf{2.} Whatever the first birthday, the second matches it with probability
$1/365$, about 0.003.

\textbf{3.} Any guess will do: most people say about 0.05 or 0.1.

\textbf{4.} \code{True} as soon as a birthday appears twice; once that has
happened, the rest of the class cannot change the answer.

\omoutput{code/grade-9/05-randomness-and-simulation/outputs/birthday.txt}{10,000
simulated classes of each size.}\label{out:g9:randomness-and-simulation:birthday}

\textbf{5.} 0.517.

\textbf{6.} From 23 pupils: 20 give 0.416, 23 give 0.517.

\textbf{7.} With 60 pupils there are so many pairs that a class without a
shared birthday is very rare: 0.995 means about 5 classes in 1,000.

\textbf{8.} For all birthdays to differ, pupil $k + 1$ must avoid the $k$ days
already taken: $365 - k$ of the 365 days. Multiplying these fractions gives the
chance that all differ; one minus it is the chance of a shared birthday.

\omoutput{code/grade-9/05-randomness-and-simulation/outputs/birthday_exact.txt}{The
exact chances.}\label{out:g9:randomness-and-simulation:exact}

\textbf{9.} 0.5073, and 23 is the smallest class above one half.

\textbf{10.} Within 0.01 for every size: 0.517 against 0.507 for 23, 0.416
against 0.411 for 20. 10,000 classes give about two decimals.

\textbf{11.} A shared birthday can be between any two pupils, not only with
you: a class of 23 has $23 \times 22 / 2 = 253$ pairs, each a chance of a
match.

\textbf{12.} \textbf{Result.} With 23 pupils, the chance of a shared birthday
is already about 0.507, more likely than not. The \omterm{def:g9:randomness-and-simulation:simulation}{simulation} of 10,000 classes
and the exact product agree within 0.01 for every size from 5 to 60.
\end{solution}
